\documentclass[10pt,a4paper]{amsart}
\usepackage{amsmath,amssymb,amsthm,graphicx}
\usepackage{verbatim}
\newenvironment{reference}{\comment}{\endcomment}
%\newenvironment{reference}{}{}
\newenvironment{slogan}{\comment}{\endcomment}
\newenvironment{history}{\comment}{\endcomment}

% For commutative diagrams we use Xy-pic
\usepackage[all]{xy}

% We use 2cell for 2-commutative diagrams.
\xyoption{2cell}
% Theorem environments.
%
\theoremstyle{plain}
\newtheorem{theorem}[subsection]{Theorem}
\newtheorem{proposition}[subsection]{Proposition}
\newtheorem{lemma}[subsection]{Lemma}

\theoremstyle{definition}
\newtheorem{definition}[subsection]{Definition}
\newtheorem{example}[subsection]{Example}
\newtheorem{exercise}[subsection]{Exercise}
\newtheorem{situation}[subsection]{Situation}

\theoremstyle{remark}
\newtheorem{remark}[subsection]{Remark}
\newtheorem{remarks}[subsection]{Remarks}

\numberwithin{equation}{subsection}

% Macros
%
\def\lim{\mathop{\mathrm{lim}}\nolimits}
\def\colim{\mathop{\mathrm{colim}}\nolimits}
\def\Spec{\mathop{\mathrm{Spec}}}
\def\Hom{\mathop{\mathrm{Hom}}\nolimits}
\def\Ext{\mathop{\mathrm{Ext}}\nolimits}
\def\SheafHom{\mathop{\mathcal{H}\!\mathit{om}}\nolimits}
\def\SheafExt{\mathop{\mathcal{E}\!\mathit{xt}}\nolimits}
\def\Sch{\mathit{Sch}}
\def\Mor{\mathop{\mathrm{Mor}}\nolimits}
\def\Ob{\mathop{\mathrm{Ob}}\nolimits}
\def\Sh{\mathop{\mathit{Sh}}\nolimits}
\def\NL{\mathop{N\!L}\nolimits}
\def\CH{\mathop{\mathrm{CH}}\nolimits}
\def\proetale{{pro\text{-}\acute{e}tale}}
\def\etale{{\acute{e}tale}}
\def\QCoh{\mathit{QCoh}}
\def\Ker{\mathop{\mathrm{Ker}}}
\def\Im{\mathop{\mathrm{Im}}}
\def\Coker{\mathop{\mathrm{Coker}}}
\def\Coim{\mathop{\mathrm{Coim}}}

% Boxtimes
%
\DeclareMathSymbol{\boxtimes}{\mathbin}{AMSa}{"02}

%
% Macros for moduli stacks/spaces
%
\def\QCohstack{\mathcal{QC}\!\mathit{oh}}
\def\Cohstack{\mathcal{C}\!\mathit{oh}}
\def\Spacesstack{\mathcal{S}\!\mathit{paces}}
\def\Quotfunctor{\mathrm{Quot}}
\def\Hilbfunctor{\mathrm{Hilb}}
\def\Curvesstack{\mathcal{C}\!\mathit{urves}}
\def\Polarizedstack{\mathcal{P}\!\mathit{olarized}}
\def\Complexesstack{\mathcal{C}\!\mathit{omplexes}}
% \Pic is the operator that assigns to X its picard group, usage \Pic(X)
% \Picardstack_{X/B} denotes the Picard stack of X over B
% \Picardfunctor_{X/B} denotes the Picard functor of X over B
\def\Pic{\mathop{\mathrm{Pic}}\nolimits}
\def\Picardstack{\mathcal{P}\!\mathit{ic}}
\def\Picardfunctor{\mathrm{Pic}}
\def\Deformationcategory{\mathcal{D}\!\mathit{ef}}
\usepackage{cleveref}
\crefname{originalreference}{original reference}{original references}
\Crefname{originalreference}{Original reference}{Original references}
\makeatletter
\newlabel{spaces-properties-section-types-properties}{{03E5}{source}{}{sourcepointer-1}{}}
\newlabel{spaces-properties-section-types-properties@cref}{{[originalreference][0][]03E5}{[1][0][]source}{}{}{}}
\makeatother
\crefname{originalreference}{original reference}{original references}
\Crefname{originalreference}{Original reference}{Original references}
\makeatletter
\newlabel{spaces-morphisms-lemma-smooth-local}{{03ZF}{source}{}{sourcepointer-2}{}}
\newlabel{spaces-morphisms-lemma-smooth-local@cref}{{[originalreference][0][]03ZF}{[1][0][]source}{}{}{}}
\makeatother
\crefname{originalreference}{original reference}{original references}
\Crefname{originalreference}{Original reference}{Original references}
\makeatletter
\newlabel{descent-lemma-reduced-local-smooth}{{034E}{source}{}{sourcepointer-3}{}}
\newlabel{descent-lemma-reduced-local-smooth@cref}{{[originalreference][0][]034E}{[1][0][]source}{}{}{}}
\makeatother
\title{Reducedness and smooth morphisms of algebraic spaces}
\author{The Stacks Project Authors}
\date{stacks-project:spaces-descent.tex:local-snapshot-20261012; GNU FDL 1.2 or later}
\begin{document}
\maketitle
\noindent\textit{Editorial note.} All delivered mathematical text is the author's own, copied verbatim from the pinned Stacks Project chapter file spaces-descent.tex; no statement or proof step was rewritten, reordered or removed.\par\medskip
\noindent\textit{Editorial note.} The counted span cites one Morphisms of Spaces lemma, one Descent lemma and one Properties of Spaces section, recorded in assembly.external\_references with canonical tags as navigation pointers.\par\medskip
\noindent\textit{Allowed standard prerequisites (editorial).}\begin{itemize}
\item Definition of reduced algebraic spaces. Stacks Project, Properties of Spaces, Section spaces-properties-section-types-properties. Equates reducedness of spaces with reducedness of charts.
\item Smoothness is \textbackslash{}'etale local on the source. Stacks Project, Morphisms of Spaces, Lemma spaces-morphisms-lemma-smooth-local. Shows \$U \textbackslash{}to V\$ is smooth.
\item Reducedness is smooth-local for schemes. Stacks Project, Descent, Lemma descent-lemma-reduced-local-smooth. Gives both directions for \$U \textbackslash{}to V\$.
\end{itemize}\medskip
\ifdefined\hypertarget\hypertarget{sourcepointer-1}{}\fi
\noindent\textit{Source reference.} Original reference 03E5: sources/stacks-chow/tags.txt, source lines 4444--4444. Cross-chapter standard prerequisite cited by the author inside the counted span; it is used as a pointer and its own proof is not part of this item.\par\medskip
\ifdefined\hypertarget\hypertarget{sourcepointer-2}{}\fi
\noindent\textit{Source reference.} Original reference 03ZF: sources/stacks-chow/tags.txt, source lines 5218--5218. Cross-chapter standard prerequisite cited by the author inside the counted span; it is used as a pointer and its own proof is not part of this item.\par\medskip
\ifdefined\hypertarget\hypertarget{sourcepointer-3}{}\fi
\noindent\textit{Source reference.} Original reference 034E: sources/stacks-chow/tags.txt, source lines 4089--4089. Cross-chapter standard prerequisite cited by the author inside the counted span; it is used as a pointer and its own proof is not part of this item.\par\medskip
\begin{lemma}
\label{lemma-reduced-local-smooth}
Let $f : X \to Y$ be a smooth morphism of algebraic spaces.
If $Y$ is reduced, then $X$ is reduced. If $f$ is surjective
and $X$ is reduced, then $Y$ is reduced.
\end{lemma}

\begin{proof}
Choose a commutative diagram
$$
\xymatrix{
U \ar[d] \ar[r] & V \ar[d] \\
X \ar[r] & Y
}
$$
where $U$ and $V$ are schemes, the vertical arrows are surjective and
\'etale, and $U \to X \times_Y V$ is surjective \'etale. Observe that $X$ is
a reduced algebraic space if and only if $U$ is a reduced scheme
by our definition of reduced algebraic spaces in
Properties of Spaces, Section \ref{spaces-properties-section-types-properties}.
Similarly for $Y$ and $V$. 
The morphism $U \to V$ is a smooth morphism of schemes, see
Morphisms of Spaces, Lemma \ref{spaces-morphisms-lemma-smooth-local}.
Since being reduced is local for the smooth topology for
schemes (Descent, Lemma \ref{descent-lemma-reduced-local-smooth})
we see that $U$ is reduced if $V$ is reduced.
On the other hand, if $X \to Y$ is surjective, then $U \to V$ is
surjective and in this case if $U$ is reduced, then $V$ is reduced.
\end{proof}

\end{document}
